\documentclass[10pt,a4paper]{article}
\usepackage{xeCJK}
\usepackage{fontspec}
\usepackage{booktabs}
\usepackage{array}
\usepackage{xcolor}
\usepackage{geometry}
\usepackage[protrusion=true]{microtype}
\geometry{margin=1.5cm}
\linespread{1.02}

\setmainfont{Baskerville}
\setCJKmainfont{Songti SC}
\newfontfamily{\starfont}{Songti SC}

\definecolor{wrongred}{RGB}{198,40,40}
\definecolor{rulegray}{RGB}{190,190,190}

\newcommand{\checkmarkbox}{\fbox{\rule{0pt}{0.75ex}\rule{0.75ex}{0pt}}}
\newcommand{\STAR}{\textcolor{wrongred}{\starfont ★}}
% \fa{题号}{答案}{解析}
\newcommand{\fa}[3]{\noindent\textbf{#1.}~\textbf{#2}\par\nobreak\vspace{0.03em}\hspace*{1.5em}{{\footnotesize #3}\par\vspace{0.16em}}}

\begin{document}
\pagestyle{plain}

\begin{center}
{\LARGE\bfseries 动名词练习填空（with-keys-1）\quad 参考答案与解析}\\[0.25em]
{\large 上海市大同中学高一英语 · 非谓语动词专项（共 50 题）}\\[0.22em]
{\small 2029届高一英语 · 解析以"主动 vs 被动、先发生 vs 后发生"为主线}
\end{center}

\vspace{0.3em}
\noindent{\small 说明：原卷未印答案（with-keys 版面仅有题面），本答案为按标准语法补拟；第 16、41、45 题老师红笔判错但作答形似常见答案，取口径见各条解析，请以讲评为准。}

\vspace{0.45em}
\fa{1}{control}{be used to do（被用来做）$\to$ 不定式，主动。}
\fa{2}{eating}{get used to（习惯于）中 to 是介词 $\to$ 动名词。}
\fa{3}{to clean}{would like to do 想要做。}
\fa{4}{achieving}{succeed in＋动名词。}
\fa{5}{being punished}{escape＋动名词；Ryan 与 punish 是被动关系 $\to$ being done。}
\fa{6}{smoking}{allow doing sth（允许某事发生）；allow sb to do sth 才接人。}
\fa{7}{to take}{be allowed to do（人不许做某事）。}
\fa{8}{to answer}{have letters to answer：不定式作后置定语（用 answer 的主动形式表被动含义）。}
\fa{9}{to deal}{be difficult to deal with：不定式，与主语是动宾关系但用主动（deal with 不拆）。}
\fa{10}{seeing}{be worth doing 主动形式表被动含义。}
\fa{11}{being made}{be worthy of＋being done（小说是被拍）。}
\fa{12}{to do}{be made to do：make 用于被动语态后要还原 to。}
\fa{13}{being operated}{insist on＋动名词；the patient 与 operate on 是被动关系 $\to$ being operated on。}
\fa{14}{to get; to come}{fail to do；get sb to do（让某人做）。}
\fa{15}{to save}{everything he could (do) 后接 to save 作目的状语。}
\fa{16}{quarrelling（存疑）}{It is any good doing 的口径应为动名词，与红笔判错有出入，请以讲评为准。}
\fa{17}{not go}{would rather not do（否定词放在 rather 之后）。}
\fa{18}{to die; give}{prefer to do rather than (to) do：than 前后同为不定式，第二个 to 可省。}
\fa{19}{telling}{regret doing：遗憾"已做过"（先于 regretted）。}
\fa{20}{to tell}{regret to do：遗憾地要做（告之坏消息）。}
\fa{21}{watch}{do nothing but do：but 前有实义动词 do，其后用不带 to 的不定式。}
\fa{22}{to do}{have no choice but to do：but 前无 do，须带 to。}
\fa{23}{having been knocked（to have been knocked 亦可）}{forget having been done：上个月"被撞过"是已发生且被动。}
\fa{24}{not to be interrupted}{what he wants is to be done；否定词 not 置于不定式前。}
\fa{25}{to leave}{be forced to do。}
\fa{26}{setting}{be devoted to 中 to 是介词 $\to$ 动名词。}
\fa{27}{to climb}{表语用不定式：goal is to climb（具体的打算）。}
\fa{28}{designing}{表语动名词；design 双写 n 再加 -ing。}
\fa{29}{to be held}{meeting 与 hold 是被动关系且 tomorrow 未发生 $\to$ to be done 作后置定语。}
\fa{30}{to be translating}{is said＋不定式；these days 表"这几天正在进行" $\to$ to be doing。}
\fa{31}{taking}{How about＋动名词；the two of us 是逻辑主语。}
\fa{32}{to speak}{needs a lot of practice to speak：不定式作状语（练习是为了说好）。}
\fa{33}{checking（to be checked 亦可）}{need doing 主动表被动；need to be checked 同义。原卷先写 to be checked 被划改，老师口径取动名词。}
\fa{34}{to practise}{sb needs to do（人需要做某事）。}
\fa{35}{to have been}{happen to have done：happened 之前"去过" $\to$ 不定式完成式。}
\fa{36}{to have been given}{glad 之前"已获得机会" $\to$ to have been done。原卷先写 to be given（被判错划改后反成 giving）。}
\fa{37}{to fall}{be seen to do：see 用于被动后还原 to。}
\fa{38}{hiding; to have been told}{It is no use doing；she 与 tell 是被动关系且"已被告知" $\to$ appears to have been told。}
\fa{39}{to be working}{pretend to be doing：when mother came in 提供"当时正在"的背景。}
\fa{40}{adding}{try doing 试着做（看看效果）；try to do 是"努力做成"。}
\fa{41}{not having written}{动名词完成式：一年没写信发生在 blamed 之前。红笔判错，或取此口径，请以讲评为准。}
\fa{42}{to improve}{help (to) do 有助于（couldn't help to do 无助于）；≠ can't help doing（情不自禁）。}
\fa{43}{to have}{stop to do 停下来去做（另一件事）；stop doing 是停止手头的事。}
\fa{44}{to be told}{only to do 表意外结果；she 与 tell 是被动关系 $\to$ only to be told。}
\fa{45}{(should) run（存疑）}{suggest＋that 从句须用 (should) do 虚拟语气；若按动名词复合结构则为 my brother's running。红笔判错，请以讲评为准。}
\fa{46}{to have been shipped}{a few days ago 已发生，goods 与 ship 是被动关系 $\to$ be believed to have been done。}
\fa{47}{to be punished}{for her to be done：it is wrong for sb to be done。原卷先写 to be punished（划改后反成 punishing）。}
\fa{48}{To be admitted}{"被吸收入党"在 write an application 之后 $\to$ 不定式被动作条件状语；动名词 Being admitted 表已发生，不合语境。}
\fa{49}{his insulting}{动名词复合结构：物主代词 his／名词所有格＋doing。}
\fa{50}{being injured}{prevent sb from＋动名词；workers 与 injure 是被动关系 $\to$ being done。原卷先写 being injured（划改后反成 injuring）。}

\vspace{0.3em}
\noindent{\color{rulegray}\rule{\textwidth}{0.5pt}}
\vspace{0.3em}
\noindent\textbf{复核记录}\hfill\small 完成重做后填写

\vspace{0.3em}
\noindent\begin{tabular}{@{}p{0.31\textwidth}@{}p{0.34\textwidth}@{}p{0.35\textwidth}@{}}
重做得分：\underline{\hspace{1.2cm}}\,/\,50
& 仍错误的题号：\underline{\hspace{3.2cm}}
& 下次复习日期：\underline{\hspace{2.0cm}} \\
\end{tabular}

\vspace{0.3em}
\noindent 本轮易错点：\checkmarkbox\ 动名词被动/完成式\quad\checkmarkbox\ 不定式被动/进行/完成式\quad\checkmarkbox\ 固定搭配\quad\checkmarkbox\ 拼写

\vspace{0.3em}
\noindent 复习状态：\checkmarkbox\ 已掌握\quad\checkmarkbox\ 需要巩固\quad\checkmarkbox\ 需重点复习

\end{document}
